\section{Metric Spaces} \label{an:sec:metric-spaces}

\begin{definition} \label{an:def:metric-space}
  A function $d: X \times X \to \R$ is called a \textbf{metric} on $X$ if for all $p, q, r \in X$,
  \begin{enumerate}[label=\textbf{(\Alph*)}]
    \item $d(p, q) \geq 0$
    \item $d(p, q) = 0 \iff p = q$
    \item $d(p, q) = d(q, p)$
    \item $d(p, r) \leq d(p, q) + d(q, r)$
  \end{enumerate}
  \textbf{(D)} is called the \textbf{triangle inequality}. The pair $(X, d)$ is called a \textbf{metric space}.

\end{definition}

\bigskip

\begin{theorem} \label{an:thm:standard-metric-on-rn}
  Suppose $n \in \N$. Define $d: \R^{n} \times \R^{n} \to \R$ as
  \[
    d(\bm{x}, \bm{y}) \coloneqq \norm*{\bm{x} - \bm{y}} = \sqrt{\sum_{k = 1}^{n} (x_{k} - y_{k})^{2}}
  \]
  for all $\bm{x}, \bm{y} \in \R^{n}$. Then $d$ is a metric on $\R^{n}$.
  In particular, $d$ is called the \textbf{standard metric} on $\R^{n}$.
\end{theorem}

\begin{proof}
  Suppose $\bm{x}, \bm{y}, \bm{z} \in \R^{n}$.
  \begin{itemize}
    \item $d(\bm{x}, \bm{y}) \geq 0$ is trivial.
    \item By \cref{an:def:inner-product-norm},
          \begin{align*}
            d(\bm{x}, \bm{y}) = 0 & \iff \norm*{\bm{x} - \bm{y}} = 0              \\
                                  & \iff \forall k \in [1, n]_{\N}, x_{k} = y_{k} \\
                                  & \iff \bm{x} = \bm{y}
          \end{align*}
    \item $d(\bm{x}, \bm{y}) = d(\bm{y}, \bm{x})$ is trivial.
    \item By \cref{an:thm:properties-of-inner-product-norme},
          \begin{align*}
            d(\bm{x}, \bm{z}) & = \norm*{\bm{x} - \bm{z}} =
            \norm*{(\bm{x} - \bm{y}) + (\bm{y} - \bm{z})}                              \\
                              & \leq \norm*{\bm{x} - \bm{y}} + \norm*{\bm{y} - \bm{z}}
            = d(\bm{x},\bm{y}) + d(\bm{y},\bm{z})
          \end{align*}
  \end{itemize}
  Therefore, $d$ is a metric on $\R^{n}$.
\end{proof}

\bigskip

\begin{definition} \label{an:def:k-cell}
  Suppose $k \in \N$. A set $E \subseteq \R^{k}$ is called a \textbf{$k$-cell} if
  \[
    E = \setbuilder*{\bm{x} \in \R^{k}}{\forall j \in [1, k]_{\N}, a_{j} \leq x_{j} \leq b_{j}}
  \]
  where
  \[
    \forall j \in [1, k]_{\N}, - \infty < a_{j} \leq b_{j} < + \infty
  \]
\end{definition}

\bigskip

\begin{definition} \label{an:def:neighborhood-open-ball}
  Suppose $(X, d)$ is a metric space. Suppose $p \in X$ and $r \in \R^{+}$.
  \begin{itemize}
    \item The \textbf{neighborhood} of $p$ with radius $r$ is defined as
          \[
            N_{r}(p) \coloneqq \setbuilder*{q \in X}{d(p, q) < r}
          \]
    \item The \textbf{deleted neighborhood} of $p$ with radius $r$ is defined as
          \[
            \widetilde{N}_{r}(p) \coloneqq N_{r}(p) \setminus \set*{p} = \setbuilder*{q \in X}{0 < d(p, q) < r}
          \]
  \end{itemize}
\end{definition}

\bigskip

\begin{definition} \label{an:def:limit-isolated-interior-points}
  Suppose $(X, d)$ is a metric space and $E \subseteq X$. Suppose $p \in X$.
  \begin{itemize}
    \item $p$ is called an \textbf{limit point} of $E$ if
          \[
            \forall r \in \R^{+}, E \cap \widetilde{N}_{r}(p) \neq \vnot
          \]
          The set of all limit points of $E$ is denoted by $E'$.
    \item $p$ is called a \textbf{isolated point} of $E$ if
          \[
            \exists r \in \R^{+}, E \cap N_{r}(p) = \set*{p}
          \]
    \item $p$ is called an \textbf{interior point} of $E$ if
          \[
            \exists r \in \R^{+}, N_{r}(p) \subseteq E
          \]
          The set of all interior points of $E$ is denoted by $E^{\circ}$.
  \end{itemize}
\end{definition}

\bigskip

\begin{definition} \label{an:def:open-closed-sets}
  Suppose $(X, d)$ is a metric space and $E \subseteq X$.
  \begin{itemize}
    \item $E$ is said to be \textbf{open} if
          \[
            \forall p \in E, p \in E^{\circ}
          \]
    \item $E$ is said to be \textbf{closed} if
          \[
            E' \subseteq E
          \]
  \end{itemize}
\end{definition}

\bigskip

\begin{definition} \label{an:def:closure}
  Suppose $(X, d)$ is a metric space and $E \subseteq X$. The \textbf{closure} $\overline{E}$ of $E$ is defined as
  \[
    \overline{E} \coloneqq E \cup E'
  \]
\end{definition}

\bigskip

\begin{definition} \label{an:def:bounded-perfect-dense-sets}
  Suppose $(X, d)$ is a metric space and $E \subseteq X$.
  \begin{itemize}
    \item $E$ is said to be \textbf{bounded} if
          \[
            \exists p \in X, \exists r \in \R^{+}, E \subseteq N_{r}(p)
          \]
    \item $E$ is said to be \textbf{perfect} if
          \[
            E = E'
          \]
    \item $E$ is said to be \textbf{dense} if
          \[
            \overline{E} = X
          \]
  \end{itemize}
\end{definition}

\bigskip

\begin{theorem} \label{an:thm:neighborhood-is-open}
  Suppose $(X, d)$ is a metric space.
  \[
    \forall p \in X, \forall r \in \R^{+}, N_{r}(p) \text{ is open}
  \]
\end{theorem}

\begin{proof}
  Suppose $q \in N_{r}(p)$. Then
  \[
    \begin{array}{c}
      d(p, q) < r \\
      \Downarrow  \\
      \veps \coloneqq r - d(p, q) > 0
    \end{array}
  \]
  Suppose $s \in N_{\veps}(q)$. Then
  \[
    \begin{array}{c}
      d(q, s) < \veps                                      \\
      \Downarrow                                           \\
      d(p, s) \leq d(p, q) + d(q, s) < d(p, q) + \veps = r \\
      \Downarrow                                           \\
      s \in N_{r}(p)                                       \\
      \Downarrow                                           \\
      N_{\veps}(q) \subseteq N_{r}(p)
    \end{array}
  \]
  Since $q$ is arbitrary, $N_{r}(p)$ is open.
\end{proof}

\bigskip

\begin{theorem} \label{an:thm:limit-point-neighborhood-contains-infinite-elements}
  Suppose $(X, d)$ is a metric space and $E \subseteq X$.
  \[
    \forall p \in E', \forall r \in \R^{+}, \abs*{E \cap N_{r}(p)} \geq \aleph_{0}
  \]
\end{theorem}
